

\chapter{Requisitos e Especificações}
\label{chapter:requirements}

\begin{introduction}
Este capítulo especifica os requisitos funcionais e não funcionais do sistema, bem como as restrições de design que condicionaram as decisões de projeto.
\end{introduction}



\section{Requisitos Funcionais}
\label{sec:requisitos-funcionais}

Os requisitos funcionais foram definidos a partir das necessidades concretas do utilizador final (investigação sísmica com recursos limitados) e das restrições técnicas impostas pelos componentes selecionados. Definem o que o sistema deve fazer para cumprir os objetivos do projeto.

\begin{itemize}
    \item \textbf{RF1 --- Movimento bi-axial:} Para simular excitações sísmicas bidimensionais, a plataforma deve mover-se de forma controlada em dois eixos horizontais independentes, designados por eixo~X (direção Norte-Sul) e eixo~Y (direção Este-Oeste).

    \item \textbf{RF2 --- Reprodução de sismo real:} O registo sísmico de El~Centro de 1940, nas suas componentes Norte-Sul e Este-Oeste, deve ser reproduzível a partir de dados de deslocamento pré-processados e codificados no firmware.

    \item \textbf{RF5 --- Frequência de atualização de 50\,Hz:} As velocidades dos motores devem ser atualizadas a 50\,Hz (passo de 20\,ms), para assegurar fidelidade temporal na reprodução dos dados sísmicos pré-processados.

    \item \textbf{RF7 --- Suporte de modelos de teste:} A plataforma superior deve ter dimensões e resistência estrutural suficientes para receber modelos reduzidos de estruturas e observar o seu comportamento dinâmico durante os ensaios.
\end{itemize}


\section{Requisitos Não Funcionais}
\label{sec:requisitos-nao-funcionais}

Para além das funcionalidades, o sistema tem de satisfazer atributos de qualidade e restrições operacionais que condicionam a forma como é construído e disponibilizado, independentemente da função específica que executa.

\begin{itemize}
    \item \textbf{RNF1 --- Custo:} O custo total do protótipo não deve ultrapassar os 500\,\euro{}, valor que o torna acessível em contextos de investigação sísmica com orçamentos limitados e significativamente mais baixo do que o das soluções comerciais equivalentes.

    \item \textbf{RNF2 --- Reprodutibilidade:} Todos os ficheiros de modelação 3D, o firmware, os scripts Python e a documentação técnica devem estar disponíveis publicamente, de modo a que qualquer investigador possa construir uma unidade idêntica sem necessidade de ferramentas de fabrico especializadas.

    \item \textbf{RNF3 --- Segurança:} Fins de curso físicos ou lógicos devem limitar o deslocamento máximo admissível da plataforma em cada eixo, para prevenir colisões destrutivas dos componentes móveis em caso de erro de controlo ou de dados corrompidos.

    \item \textbf{RNF4 --- Facilidade de montagem:} Qualquer pessoa com conhecimentos técnicos básicos em eletrónica e impressão 3D deve conseguir montar o sistema, sem recurso a maquinagem especializada, equipamento de soldadura complexo ou ferramentas industriais.

    \item \textbf{RNF5 --- Compatibilidade com ecossistema \textit{open-source}:} Tanto o firmware como os scripts de pré-processamento devem assentar em ferramentas \textit{open-source}, nomeadamente o Arduino IDE e a biblioteca AccelStepper, para assegurar compatibilidade com o ecossistema Arduino/ESP32 e facilitar adaptações futuras por parte de terceiros.
\end{itemize}


\section{Restrições de Design}
\label{sec:restricoes-design}

As restrições de design delimitam o espaço tecnológico e económico do projeto. Algumas resultam de opções deliberadas orientadas para o baixo custo e a reprodutibilidade; outras decorrem das limitações inerentes aos componentes selecionados.

\begin{itemize}
    \item \textbf{RD1 --- Componentes de baixo custo e disponibilidade comercial:} Recorre-se exclusivamente a componentes disponíveis em fornecedores de eletrónica genérica, nomeadamente motores NEMA17, \textit{drivers} DRV8825, microcontrolador ESP32 e CNC Shield. A translação linear baseia-se em fusos trapezoidais com \textit{lead} de 8\,mm por revolução, escolha que equilibra custo, precisão e capacidade de carga sem recorrer a componentes proprietários ou de difícil aquisição em mercados internacionais.

    \item \textbf{RD2 --- Estrutura mecânica em impressão 3D:} A estrutura mecânica é fabricada exclusivamente por impressão 3D, sem recurso a maquinagem CNC ou outros processos subtrativos, de forma a garantir a reprodutibilidade por qualquer utilizador com acesso a uma impressora de gama de entrada.

    \item \textbf{RD4 --- Microcontrolador ESP32 com CNC Shield:} Na camada de controlo, utiliza-se o ESP32 em conjunto com uma CNC Shield compatível, que providencia a interface normalizada entre o microcontrolador e os \textit{drivers} dos motores e simplifica consideravelmente a cablagem.

    \item \textbf{RD5 --- Linguagens de programação:} O firmware é desenvolvido em C++ com suporte para Arduino (ESP32) e o pré-processamento dos dados sísmicos em Python, por serem as linguagens mais difundidas na comunidade académica e de \textit{makers} e por disporem de um vasto ecossistema de bibliotecas \textit{open-source}.
\end{itemize}
